\chapter{Beyond the wall: operating-point tomography}
\label{ch:tomography}

\begin{quote}\itshape
This chapter in brief: chapter 8 ended with ``buy sensors.'' This chapter finds two cheaper
exits. First, the \emph{schedule} of routine snapshot reports is itself an observability
design variable --- one added report condition buys $+79\,\%$ separability on the HPT
ambiguity with zero hardware. Second, a learner that fuses per-flight \emph{physics
projections} over time more than triples classical multi-point isolation on identical data.
Both claims survived a second pre-registered freeze; one companion hypothesis did not, and
that story is told too.
\end{quote}

\section{The idea, and what free variation turns out to be worth}
\label{sec:tomo-idea}

The influence matrix is not one matrix: $\mathbf{H}(u)$ changes with the flight condition.
A window of ordinary snapshots at scattered conditions --- hot day, derate, another flight
level --- is therefore a set of \emph{different projections of the same health state}, and
fusing them is a tomography problem. Classical multi-operating-point analysis exploits this
with deliberately scheduled test points and a constant-health assumption; the question here
is what the \emph{in-service} version is worth.

The physics audit (\code{scripts/f7_observability.py}, \cref{fig:f7-obs}) gave a two-sided
answer. Encouraging: today's dual-report structure (one takeoff, one cruise snapshot per
flight) already restores full algebraic rank --- 10/10 from two flights --- because the two
report families are far apart in condition space. Sobering: the \emph{free} ambient and
power scatter around those two conditions saturates almost immediately; the critical
signature angles improve only $\approx 1.4\times$ over a single point, far short of what
deliberately spread conditions achieve. Free variation is real but thin.

\begin{figure}[htbp]
  \centering
  \includegraphics[width=0.92\textwidth]{f7_observability}
  \caption{Observability under realistic in-service scatter (200 draws per window size).
  Left: identifiable rank recovers immediately. Right: signature angles saturate ---
  free scatter is too narrow to keep paying. Auto-generated.}
  \label{fig:f7-obs}
\end{figure}

\section{The report schedule as a design variable}
\label{sec:tomo-schedule}

If free scatter saturates but designed spread works, the actionable question becomes:
\emph{which report condition should be added to the monitoring schedule?} Treating the
ACARS report schedule as an observability design problem --- to our knowledge a new framing
--- and greedily evaluating candidate additions on the influence-matrix grid gives
\cref{tab:schedule}: a single periodic \textbf{stabilized low-power cruise report} nearly
doubles the separability of the HPT efficiency-vs-flow ambiguity ($6.2^{\circ} \to
11.0^{\circ}$, $+79\,\%$); the best pair of additions reaches $12.7^{\circ}$. The cost is a
procedure bulletin, not avionics. Equally important is what no schedule fixes:
$\eta_{HPC}\!\sim\!\eta_{HPT}$ stays below $1.8^{\circ}$ everywhere in the envelope ---
that pair is \emph{fundamental} for cockpit sensors, and the sensor-purchase conclusion of
chapter 8 stands for it alone.

\begin{table}[htbp]
  \centering\small
  \caption{Marginal diagnosability value of adding one report condition to today's
  schedule (cockpit sensors). Auto-generated from the ICM grid.}
  \label{tab:schedule}
  % AUTOGENERATED by scripts/make_report_assets.py — do NOT edit by hand
\begin{tabular}{lrr}
\toprule
Report schedule & $\eta_{HPC}\!\sim\!\eta_{HPT}$ [deg] & $\eta_{HPT}\!\sim\!\Gamma_{HPT}$ [deg] \\
\midrule
  today's schedule (takeoff + cruise) & 1.33 & 6.16 \\
  + low-power stabilized cruise & 1.64 & 11.03 \\
  + climb report & 1.36 & 7.33 \\
  + FL390 cruise & 1.41 & 8.04 \\
  + hot-day takeoff & 1.36 & 7.52 \\
  + best pair (low-power cruise, hot-day takeoff) & 1.71 & 12.65 \\
\bottomrule
\end{tabular}

\end{table}

\section{Learned MOPA: fusing physics projections over time}
\label{sec:tomo-learner}

The classical way to fuse a window of $T$ flights is to stack every snapshot's linear
observation into one tall system and solve it as a single regularized least squares:
\begin{equation}
  \begin{bmatrix} \Delta\mathbf{z}_1 \\ \vdots \\ \Delta\mathbf{z}_T \end{bmatrix}
  =
  \begin{bmatrix} \mathbf{H}(u_1) \\ \vdots \\ \mathbf{H}(u_T) \end{bmatrix} \mathbf{x}
  + \mathbf{v},
  \qquad
  \hat{\mathbf{x}} = \bigl(\mathbf{H}_s^{\!\top}\mathbf{H}_s + \lambda \mathbf{I}\bigr)^{-1}
  \mathbf{H}_s^{\!\top} \Delta\mathbf{z}_s,
  \label{eq:stacked-wls}
\end{equation}
predicting the fault as the largest-magnitude entry of $\hat{\mathbf{x}}$ (with a
no-fault threshold). The single $\mathbf{x}$ on the right-hand side is the method's load-
bearing assumption --- \emph{one constant health state across the whole window} --- valid
on a test bed, false in service, where health drifts and washes intervene. The learner
replaces that assumption: each flight gets its own per-block projection
$\hat{\mathbf{x}}_t$ (the one-flight version of \cref{eq:stacked-wls}), and the GRU
(\cref{eq:gru}) reads the \emph{sequence} of projections, free to learn that a slow common
drift is aging while a coherent recent step is a fault. The architecture
that works is neither classical stacking nor a raw neural network: a raw GRU fed deviation
sequences failed outright (12\,\% isolation --- it must learn the entire inverse physics
from a few hundred windows), while a GRU fed the \emph{per-flight regularized WLS
projections} $\hat{\mathbf{x}}_t$ --- physics inversion first, learned temporal fusion
second --- is data-efficient exactly where the classical method is brittle. Development
also surfaced two reusable engineering facts, both now in the code: block-averaging inside
the window (not subsampling) buys the full $\sqrt{n}$ noise reduction both methods deserve,
and the drift-degradation signature of constant-health stacking appeared in development
precisely as predicted.

\section{Second freeze, second verdicts}
\label{sec:tomo-verdicts}

The confirmatory protocol was frozen a second time (\code{docs/prereg-v2.md}, tag
\code{prereg-v2}: thresholds, pinned methods, development on a held-out slice of
\emph{training} engines so validation stayed clean for conformal calibration and the test
episodes stayed unseen). One pass over the 23 test episodes (\cref{tab:f7-verdicts}):

\begin{table}[htbp]
  \centering\small
  \caption{F7 pre-registered verdicts (single confirmatory pass). Auto-generated.}
  \label{tab:f7-verdicts}
  % AUTOGENERATED by scripts/make_report_assets.py — do NOT edit by hand
\begin{tabular}{>{\raggedright\arraybackslash}p{0.20\textwidth}>{\raggedright\arraybackslash}p{0.17\textwidth}>{\raggedright\arraybackslash}p{0.51\textwidth}}
\toprule
Hypothesis & Verdict & Confirmatory evidence \\
\midrule
  H7.2$'$ fusion beats stacking & CONFIRMED & GRU 0.65 vs WLS 0.17 ($+$48 pp); McNemar $p \leq 0.0096$ \\
  H7.3 drift robustness & REFUTED (as frozen) & WLS degraded 17 pp (needed $\geq 20$); GRU 13 pp (needed $\leq 10$); pattern present \\
  H7.4$'$ calibrated ambiguity & CONFIRMED & coverage 0.96; fundamental sets 3.0 $>$ other 2.0 \\
\bottomrule
\end{tabular}

\end{table}

\begin{itemize}
  \item \textbf{H7.2$'$ confirmed, by a wide margin.} The learned fusion isolates at
    0.65 against classical stacking's 0.17 on identical windows --- and the McNemar bound (\cref{eq:mcnemar})
    holds under \emph{any} pairing of outcomes, so no replay of the test pass was needed.
    Recall from chapter 8 that the best single-snapshot method managed 0.31: temporal
    fusion of physics projections roughly doubles that too.
  \item \textbf{H7.4$'$ confirmed.} Conformal isolation sets cover at 0.96 with median
    sizes 2--3.5, and the fundamental pair receives systematically larger sets than the
    rest: ambiguity is now a calibrated, physics-explained deliverable.
  \item \textbf{H7.3 refuted as frozen --- discipline over narrative.} The drift-robustness
    thresholds (classical must degrade $\geq 20$ pp, learner $\leq 10$ pp) were calibrated
    on a 16-episode development sample and the test came out softer on both sides (17 and
    13 pp). The qualitative pattern is exactly as hypothesized --- stacking collapses to
    0.00 on long windows while the learner holds 0.52 --- but the registered claim is the
    thresholds, and the thresholds failed. The lesson is itself transferable: freezing
    quantitative margins from small development samples is how honest protocols get
    refuted, and reporting that beats renegotiating it.
\end{itemize}

\section{What F7 adds to the adoption picture}

Chapter 8's checklist gains two rows: \textbf{(6) redesign the report schedule before
buying anything} --- one procedural report condition recovers most of what deliberate
condition spread can offer; and \textbf{(7) when fusing windows, learn the fusion but keep
the physics inversion} --- projections-then-GRU tripled classical stacking where raw
learning failed. And one row hardens: for the truly fundamental ambiguities, neither
schedules nor learning substitute for instrumentation.
